% Shrnutí okruhu Přírodou inspirované počítání — jedna otázka = jedna strana.
% Kondenzát z 02-prirodou-inspirovane-pocitani.tex; slouží k opakování těsně před zkouškou.
% Struktura strany: znění otázky → osnova odpovědi → fakta v blocích → pasti.
\documentclass[10pt,a4paper]{article}
\usepackage[utf8]{inputenc}
\usepackage[T1]{fontenc}
\usepackage{lmodern}          % vektorové fonty — nutné pro expansion v microtype
\usepackage[czech]{babel}
\usepackage[margin=1.25cm,top=1.15cm,bottom=1.25cm]{geometry}
\usepackage{amsmath,amssymb}
\usepackage{parskip}
\usepackage{enumitem}
\setlist{nosep,leftmargin=1.1em,labelsep=0.35em,topsep=1pt}
\usepackage{multicol}
\setlength{\columnsep}{5mm}
\usepackage{booktabs}
\usepackage{array}
\usepackage{xcolor}
\definecolor{linkblue}{RGB}{20,60,150}
\definecolor{defbg}{RGB}{240,244,252}
\definecolor{defframe}{RGB}{100,130,190}
\definecolor{headcol}{RGB}{25,70,140}
\definecolor{markcol}{RGB}{176,84,0}
\usepackage[most]{tcolorbox}
\usepackage{microtype}
\usepackage[colorlinks=true,linkcolor=linkblue,urlcolor=linkblue]{hyperref}
\usepackage{fancyhdr}

\pagestyle{fancy}
\fancyhf{}
\renewcommand{\headrulewidth}{0.4pt}
\fancyhead[L]{\footnotesize\textsc{Přírodou inspirované počítání --- shrnutí ke SZZ}}
\fancyhead[R]{\footnotesize\thepage}

% Nadpis otázky: číslo + doslovné znění věty oficiálního okruhu.
\newtcolorbox{qbox}[1]{colback=defbg,colframe=defframe,boxrule=0.7pt,arc=1.2mm,
  left=2.5mm,right=2.5mm,top=1.2mm,bottom=1.2mm,#1}
\newcommand{\qpage}[2]{%
  \begin{qbox}{}
  \textbf{\large Otázka #1.}\quad\textbf{#2}
  \end{qbox}
  \vspace{-0.3em}}
\newcommand{\hh}[1]{\par\smallskip\noindent{\bfseries\color{headcol}#1}\par\nobreak\vspace{-0.2em}}
\newcommand{\ph}[1]{\par\smallskip\noindent{\bfseries\color{markcol}#1}\par\nobreak\vspace{-0.2em}}
\newcommand{\ok}[1]{{\bfseries\color{markcol}#1)}\,}   % krok osnovy odpovědi
\newcommand{\tm}[1]{\textbf{#1}}
\newcommand{\en}[1]{\emph{#1}}
\newcommand{\ra}{$\rightarrow$}

\begin{document}
\raggedcolumns   % nestahovat/neroztahovat sloupce svisle

%=====================================================================
% Strana 1 — rozcestník
%=====================================================================
\thispagestyle{empty}
\begin{center}
{\LARGE\bfseries Přírodou inspirované počítání}\\[0.35em]
{\large Shrnutí okruhu ke státní závěrečné zkoušce --- jedna otázka na jednu stranu}\\[0.5em]
{\small SZZ Umělá inteligence, zaměření Inteligentní agenti, MFF UK \quad$\cdot$\quad srpen 2026}
\end{center}

\begin{qbox}{}
\small\raggedright
\textbf{Oficiální znění okruhu} (zdroj pravdy, 2025/2026):
{\footnotesize\url{https://www.mff.cuni.cz/cs/studenti/bc-a-mgr-studium/studijni-plany/2025-2026/informatika/mgr/umela-inteligence}}
\end{qbox}

\footnotesize
Šest vět okruhu, šest stran. Podrobný výklad v~\path{02-prirodou-inspirovane-pocitani.pdf}
(anglicky \path{ENG/02-nature-inspired-computing.pdf}); zdrojové slidy \path{materialy/}
(EVA~I --- NAIL025, EVA~II --- NAIL086, Neruda).

\begin{enumerate}[leftmargin=2em,itemsep=1pt]
\item Genetické algoritmy, genetické a~evoluční programování. \hfill \textit{str.~2}
\item Teorie schémat, pravděpodobnostní modely jednoduchého genetického algoritmu.
  \hfill \textit{str.~3}
\item Evoluční strategie, diferenciální evoluce, koevoluce, otevřená evoluce. \hfill \textit{str.~4}
\item Rojové optimalizační algoritmy. \hfill \textit{str.~5}
\item Memetické algoritmy, hill climbing, simulované žíhání. \hfill \textit{str.~6}
\item Aplikace evolučních algoritmů (evoluce expertních systémů, neuroevoluce, řešení
  kombinatorických úloh, vícekriteriální optimalizace). \hfill \textit{str.~7}
\end{enumerate}

\begin{qbox}{}
\footnotesize\raggedright
\textbf{Jednotící pohled --- tímhle lze začít kteroukoli odpověď.}
Všechny probírané metody jsou \emph{stochastické black-box hledání}: k~dispozici jsou jen
hodnoty $f$ ve zvolených bodech (žádný gradient), cenou je počet vyhodnocení.
Algoritmus drží \emph{rozdělení} $p_\theta$ nad prostorem řešení a~opakuje
\textbf{vzorkuj \ra{} vyhodnoť \ra{} aktualizuj $\theta$}; $\theta$ je jediná informace, která
přechází mezi iteracemi (co se naučil $+$ kde hledat dál). Liší se jen tím, co je $\theta$:
\textbf{EA} rodičovská populace $\cdot$ \textbf{ES} gaussián (střed, $\sigma$) $\cdot$
\textbf{CMA-ES} $(\vec m,\sigma,\mathbf C,\vec p_\sigma,\vec p_c)$ $\cdot$
\textbf{EDA} parametry modelu (marginály, bayesovská síť) $\cdot$
\textbf{PSO} pozice, rychlosti a~paměť $\cdot$ \textbf{ACO} feromonová matice $\cdot$
\textbf{SA} aktuální bod $+$ teplota $\cdot$ \textbf{tabu} aktuální bod $+$ zakázané tahy.
\end{qbox}

\vspace{0.2em}
\begin{center}
\footnotesize
\begin{tabular}{@{}p{1.35cm}p{3.0cm}p{3.9cm}p{2.9cm}p{4.4cm}@{}}
\toprule
\textbf{alg.} & \textbf{reprezentace} & \textbf{variace} & \textbf{selekce} &
\textbf{co si zapamatovat} \\
\midrule
SGA & binární řetězec & 1-bod.\ křížení, bit-flip & ruleta, generační &
  teorie schémat, BBH \\
ES & $\mathbb{R}^n$ $+$ $\vec\sigma$ & gaussovská mutace, rekombinace &
  náhodní rodiče, $(\mu,\lambda)$/$(\mu{+}\lambda)$ & pravidlo 1/5, samoadaptace \\
CMA-ES & $\mathbb{R}^n$, model $\mathcal N(\vec m,\sigma^2\mathbf C)$ & vzorkování z~rozdělení &
  pořadová, $\mu$ nejlepších & evoluční cesty, rank-1/rank-$\mu$ \\
DE & $\mathbb{R}^n$ & $\vec x_a + F(\vec x_b-\vec x_c)$, binom.\ křížení &
  souboj rodič--potomek & měřítko kroku z~populace \\
EP & doménová (automat) & jen mutace & turnaj rodiče $+$ potomci & genotyp $=$ fenotyp \\
GP & syntaktický strom & výměna podstromů, mutace & turnaj & bloat, ADF \\
PSO & $\mathbb{R}^n$ $+$ rychlost & pohybové rovnice & žádná (jen paměť) &
  $w$, $c_1$, $c_2$, gbest/lbest \\
ACO & konstrukce cesty & feromon $+$ heuristika & implicitní &
  $\tau^\alpha\eta^\beta$, odpařování \\
SA & jedno řešení & soused & Metropolis & rozvrh chlazení \\
MA & libovolná & EA operátory $+$ lokální hledání & libovolná & Lamarck vs.\ Baldwin \\
NSGA-II & libovolná & SBX $+$ polynom.\ mutace & fronty $+$ crowding &
  elitismus, crowding distance \\
NEAT & seznam hran & mutace struktury, křížení dle ID & sdílená fitness v~druhu &
  inovační čísla, speciace \\
\bottomrule
\end{tabular}
\end{center}

\begin{multicols}{2}
\footnotesize

\hh{Průřezová témata --- ptají se napříč otázkami}
\begin{itemize}
\item \tm{Explorace vs.\ exploatace} a~kde ji každý algoritmus řídí: GA (křížení vs.\ selekční
  tlak), ES ($\sigma$, pravidlo 1/5), SA (teplota), PSO ($w$, topologie), ACO ($q_0$,
  odpařování), tabu (tabu list), novelty search (jen explorace).
\item \tm{Udržení diverzity}: fitness sharing, crowding, speciace, ostrovní model, nenáhodné
  páření, restart, hypermutace.
\item \tm{Adaptace parametrů}: ladění vs.\ řízení; pevné / adaptivní / samoadaptivní.
\item \tm{Práce s~omezeními}: penalizace, oprava, dekodér, uzavřené operátory, vícekriteriální
  formulace.
\item \tm{Lamarckismus vs.\ baldwinismus}: memetické algoritmy, oprava jedinců, neuroevoluce
  s~doučováním vah.
\item \tm{Teorie}: věta o~schématech a~její kritika, Voseho model, Markovovy řetězce
  a~konvergence elitistického GA, No Free Lunch, konvergence SA.
\end{itemize}

\hh{Co se chce spočítat u~tabule}
\begin{itemize}
\item Ruletová selekce pro danou populaci; jeden krok DE s~čísly.
\item PMX a~OX na devítiprvkovém příkladu; inverze (2-opt) na permutaci.
\item Věta o~schématech: $o(S)$, $d(S)$, celé odvození nerovnosti.
\item Crowding distance v~NSGA-II; dominance a~Pareto fronta z~tabulky kritérií.
\item Metropolisovo kritérium pro konkrétní $\Delta$ a~$T$; samoadaptace $\sigma$.
\end{itemize}

\columnbreak

\hh{Zařazení oboru --- na úvodní otázku}
UI se dělí na \uv{tvrdé} logicko-formální metody a~\uv{měkké} přírodou inspirované.
\tm{Computational intelligence} (dříve soft computing) $=$ neuronové sítě $+$ fuzzy $+$
\tm{evoluční výpočty (EC)}. EC zastřešuje biologicky motivované metaheuristiky;
\tm{EVA} je jeho podmnožina stavějící na principech evoluce (populace, stochastičnost, fitness,
mutace, křížení). Mimo biologickou inspiraci stojí tabu, scatter a~novelty search.

\ph{Nejčastější chyby}
\begin{itemize}
\item Zaměnit \emph{adaptivní} (pravidlo 1/5) a~\emph{samoadaptivní} ($\sigma$ v~genomu) řízení.
\item U~samoadaptace mutovat $x$ dřív než $\sigma$ --- je to naopak (nový krok se musí hned
  otestovat).
\item Tvrdit, že věta o~schématech dokazuje konvergenci GA --- je to nerovnost na \emph{jeden}
  krok a~o~\emph{střední hodnotě}.
\item Zaměnit $(\mu,\lambda)$ a~$(\mu+\lambda)$, resp.\ jejich elitismus.
\item Považovat No Free Lunch za argument proti EA --- je to argument \emph{pro} doménovou
  specializaci reprezentace a~operátorů.
\item Říct, že lineární skalarizace vícekriteriální úlohy najde celou Pareto frontu ---
  nenajde nekonvexní části.
\item Zaměnit Pareto \emph{množinu} (v~prostoru řešení) a~Pareto \emph{frontu} (v~prostoru
  kritérií).
\end{itemize}

\end{multicols}

\newpage

%=====================================================================
\qpage{1}{Genetické algoritmy, genetické a~evoluční programování}
%=====================================================================
\begin{multicols}{2}
\footnotesize

\hh{Osnova odpovědi}
\ok{1} obecné schéma EA a~role variace vs.\ selekce \ok{2} selekce \ok{3} SGA a~binární kódování
\ok{4} reprezentace $+$ operátory $=$ fitness krajina \ok{5} GP a~jeho patologie \ok{6} varianty
GP \ok{7} EP a~srovnání čtyř historických větví.

\hh{Obecné schéma EA}
Populační stochastické prohledávání; cyklus \emph{inicializace \ra{} ohodnocení \ra{}
rodičovská selekce \ra{} rekombinace \ra{} mutace \ra{} ohodnocení \ra{} environmentální
selekce}. \tm{Variace} (křížení, mutace) diverzitu \emph{tvoří}, \tm{selekce} ji \emph{odebírá}
a~směřuje hledání; návrh EA $=$ hledání rovnováhy \tm{explorace/exploatace}.
\tm{No Free Lunch} (Wolpert \& Macready 1995--97): průměrně přes všechny úlohy si všechny
algoritmy vedou stejně \ra{} smysl má \emph{doménová specializace} reprezentace a~operátorů.

\hh{Selekce}
\tm{Ruleta} ($p_i=f_i/\sum_j f_j$; vysoký rozptyl, citlivá na škálování fitness),
\tm{SUS} (jedno náhodné místo a~dál posuny po $1/n$ --- stejné střední hodnoty, nižší rozptyl,
\uv{spravedlivější ruleta}), \tm{turnaj} (tlak řídí velikost $k$, typicky 2--5; invariantní
k~\emph{monotónním} transformacím fitness a~nepotřebuje explicitní fitness --- stačí porovnání,
proto je jediný použitelný, když se hraje hra nebo běží simulace),
\tm{pořadová} (konstantní tlak), \tm{Boltzmannova} (tlak roste s~časem).
\tm{Elitismus} zaručuje monotonii nejlepší fitness. \tm{Generační} vs.\ \tm{steady-state} model
se liší podílem nahrazované populace.
\tm{Škálování fitness} (lineární, sigma truncation, mocninné, pořadové, Boltzmannovo) opravuje
dvě slabiny ruletové selekce: dominanci \uv{supermanů} na začátku a~ztrátu tlaku na konci běhu.

\hh{SGA a~binární kódování}
\tm{SGA} $=$ binární řetězce $+$ ruleta $+$ 1-bodové křížení $+$ bit-flip mutace $+$ generační
náhrada; parametry $n$, $p_C\approx 0{,}7$, $p_M\approx 1/m$.
Reálná čísla binárně: $x=a+z(b-a)/(2^k-1)$; \tm{Hammingovy útesy} řeší \tm{Grayův kód}.
Hollandův argument pro binární abecedu (nejvíc schémat na bit) dnes \emph{neobstojí} --- binární
kódování je pro řadu úloh nepřirozené.
Křížení: \tm{1-bodové} (rozbije schéma s~pravd.\ $d(S)/(m-1)$, poziční bias),
\tm{2-bodové} (bez okrajového biasu), \tm{uniformní} (nejvíc disruptivní, bez pozičního biasu).
Mutace je v~GA \emph{sekundární} --- brání nevratné ztrátě alel; \tm{inverze} je historický
operátor pro evoluci uspořádání genů.

\hh{Reprezentace a~operátory}
\emph{Reprezentace $+$ operátory $=$ fitness krajina.}
\tm{Celočíselná}: mutace \emph{nezaujatá} (nová hodnota z~celé domény) nebo \emph{zaujatá}
(\en{creep}, normální posun) --- zaujatá dává smysl jen u~ordinálních atributů, ne
u~nominálních.
\tm{Reálná}: gaussovská mutace $x'=x+\sigma N(0,1)$; křížení \emph{strukturní} (1-bodové,
uniformní) nebo \emph{aritmetické} $\alpha\vec x+(1-\alpha)\vec y$ --- to však
\emph{zmenšuje rozptyl} populace, proto BLX-$\alpha$ a~SBX.
\tm{Permutace}: PMX (pozice $+$ mapování), OX (relativní pořadí), CX (absolutní pozice pomocí
cyklů), ER (hrany --- nejlepší pro TSP); mutace swap, insert, inverze (2-opt), scramble.
\emph{Umět předvést PMX a~OX na devítiprvkovém příkladu.}

\hh{Stromové GP (Koza 1989)}
Terminály $+$ neterminály, požadavky \tm{uzavřenosti} (každý neterminál zpracuje každou hodnotu)
a~\tm{dostatečnosti} (řešení je v~prostoru vyjádřitelné);
inicializace \en{grow} / \en{full} / \en{ramped half-and-half}; křížení $=$ výměna podstromů
(neterminály volit s~vyšší pravděpodobností, jinak se mění jen listy), víc typů mutací včetně
mutace konstant; fitness $=$ spuštění programu na datech.
\tm{Bloat} $=$ růst velikosti bez zlepšení fitness (\en{introny}); tři vysvětlení ---
\en{replication accuracy}, \en{removal bias}, \en{crossover bias}. Obrana: limity velikosti,
penalizace (\en{parsimony pressure}), anti-bloat operátory (hoist, shrink, size-fair),
vícekriteriální přístup.
\tm{ADF} $=$ evolvované podprogramy (modularita), operátory pracují na větvích odděleně;
typované a~gramatické GP omezují prostor na smysluplné programy.

\hh{Varianty GP}
\tm{Lineární GP} (byte kód; jednodušší operátory a~rychlá emulace, ale mutace snadno vyrobí
nesmysl; oblíbené v~a-life a~pro boty ve hrách). \tm{Grafové GP} (DAG --- obvody, automaty,
neuronové sítě). \tm{Cartesian GP} (lineární genotyp \ra{} 2D mřížka uzlů, bodová mutace,
$(1+4)$-ES; část uzlů zůstává neaktivní). \tm{Vývojové GP} (strom je program, který staví graf).
\tm{Gramatická evoluce}: lineární genom celých čísel, pravidlo gramatiky se vybírá
\emph{modulem}, derivační strom \ra{} fenotyp; robustní kódování a~standardní operátory GA, ale
nelokální mapování genotyp--fenotyp.

\hh{Evoluční programování a~poučení}
\tm{EP}: konečné automaty, \emph{genotyp $=$ fenotyp} (a proto \uv{křížení je zlo}), jen mutace,
turnajová selekce z~rodičů i~potomků; příklad predikce prvočísel.
Umět vyjmenovat čtyři historické větve \tm{GA / ES / EP / GP} a~jejich charakteristické volby.
\tm{Headless chicken test}: porovnej skutečné křížení s~křížením s~náhodně vygenerovaným
partnerem. Vyhraje-li náhodné, křížení nekombinuje stavební bloky, ale funguje jen jako
\emph{makromutace} --- špatné kódování, špatný operátor, nebo úloha bez stavebních bloků.

\end{multicols}
\newpage

%=====================================================================
\qpage{2}{Teorie schémat, pravděpodobnostní modely jednoduchého genetického algoritmu}
%=====================================================================
\begin{multicols}{2}
\footnotesize

\hh{Osnova odpovědi}
\ok{1} schéma a~jeho tři charakteristiky \ok{2} věta o~schématech a~její odvození ve třech
krocích \ok{3} BBH, implicitní paralelismus, bandita \ok{4} kritika \ok{5} Voseho model
(nekonečná populace) \ok{6} Markovův řetězec (konečná populace) \ok{7} EDA jako praktický
důsledek.

\hh{Schéma}
Slovo nad $\{0,1,*\}$. \tm{Řád} $o(S)$ $=$ počet pevných pozic; \tm{definující délka} $d(S)$ $=$
vzdálenost první a~poslední pevné pozice; $F(S)$ $=$ průměrná fitness reprezentantů
\emph{v~dané populaci} --- tedy veličina závislá na populaci, ne na úloze.
Počty: schémat délky $m$ je $3^m$, jeden jedinec je reprezentantem $2^m$ z~nich, schéma s~$r$
hvězdičkami pokrývá $2^r$ jedinců a~populace $n$ jedinců obsahuje mezi $2^m$ a~$n\cdot2^m$
schémat.

\hh{Věta o~schématech (TST)}
$$\mathbb{E}[C(S,t{+}1)] \ge C(S,t)\,\frac{F(S)}{\bar f}
\Bigl[1 - p_C \frac{d(S)}{m-1} - p_M\, o(S)\Bigr]$$
Odvození ve třech krocích: \emph{selekce} \ra{} faktor $F(S)/\bar f$ (nadprůměrná schémata
rostou exponenciálně), \emph{křížení} \ra{} schéma se rozbije s~pravděpodobností nejvýš
$d(S)/(m-1)$, \emph{mutace} \ra{} přežije s~pravděpodobností $(1-p_M)^{o(S)}\approx 1-p_M o(S)$.
Slovně (Holland): \emph{krátká, nízkého řádu a~nadprůměrná} schémata rostou v~populaci
exponenciálně.

\hh{Příklad na dosazení ($m=8$, $p_C=0{,}7$, $p_M=0{,}01$, $F(S)/\bar f=1{,}5$)}
$S=\mathtt{1{*}{*}0{*}{*}{*}1}$: $o(S)=3$, $d(S)=7$ \ra{}
$1-0{,}7\cdot\tfrac{7}{7}-0{,}01\cdot3=0{,}27$, tedy $1{,}5\cdot0{,}27=0{,}41$ --- schéma
\emph{ubývá}, přestože je nadprůměrné.
$S'=\mathtt{110{*}{*}{*}{*}{*}}$: $o=3$, $d=2$ \ra{}
$1-0{,}7\cdot\tfrac{2}{7}-0{,}03=0{,}77$, tedy $1{,}5\cdot0{,}77=1{,}16$ --- \emph{roste}.
Přesně to je obsah slov \uv{krátká a~nízkého řádu}: rozhoduje $d(S)$, ne jen fitness.

\hh{BBH a~implicitní paralelismus}
\tm{Hypotéza stavebních bloků}: GA skládá řešení z~krátkých, nízkého řádu a~nadprůměrných
schémat. Důsledky: záleží na \emph{kódování} (příbuzné geny blízko sebe) i~na \emph{velikosti
populace}, škodí předčasná konvergence.
\tm{Implicitní paralelismus}: populace $n$ jedinců implicitně zpracovává $2^m$ až $n\cdot2^m$
schémat, ale exponenciálně roste jen řádově $n^3$ z~nich (\uv{jediný případ, kdy je
kombinatorická exploze na naší straně}).

\hh{Bandita}
$N$ mincí, dvouruký bandita s~výplatami o~středech $\mu_1,\mu_2$ a~rozptylech
$\sigma_1,\sigma_2$. Analytické řešení: \emph{exponenciálně} více pokusů dosud vítězné páce,
$N-n^*=O(e^{c\,n^*})$. GA alokuje sloty v~populaci stejně, a~řeší tak dilema
explorace/exploatace optimálně.
Přesněji: schémata řádu $k$ soupeří o~týchž $k$ pevných pozic, hrají tedy
\tm{$2^k$-rukého banditu} (ne jednu hru s~$3^m$ pákami).

\columnbreak

\hh{Kritika teorie schémat}
\begin{itemize}
\item Nerovnost platí jen na \emph{jeden} krok a~o~\emph{střední hodnotě} --- nelze iterovat,
  $F(S)$ i~$\bar f$ se mění; v~malé populaci navíc vznikají výběrové chyby (drift).
\item Ramena bandity jsou nezávislá, schémata \emph{nejsou} --- SGA je nevzorkuje nezávisle,
  a~proto \emph{špatně odhaduje jejich fitness}. Kanonický protipříklad: $f(x)=2$ pro
  $111*\dots*$, $1$ pro $0*\dots*$, jinak $0$. Skutečně je $F(1*\dots*)=1/2 < F(0*\dots*)=1$,
  ale SGA odhadne $F(1*\dots*)\approx2$, protože v~populaci převládnou zástupci $111*\dots*$.
\item \tm{Statická BBH} (Grefenstette 1991): lidé očekávají konvergenci ke schématům se
  skutečnou průměrnou fitness, GA však konverguje k~fitness \emph{pozorované v~populacích}.
  Odtud \tm{kolaterální konvergence} (po konvergenci prefixu se ostatní schémata vzorkují
  zaujatě) a~problém \tm{velkého rozptylu fitness} uvnitř schématu.
\item \tm{Deceptivní úlohy} (stavební bloky vedou pryč od optima); TST ignoruje
  \emph{konstruktivní} efekty křížení a~mutace, počítá jen s~destruktivními.
\end{itemize}

\hh{Voseho model (nekonečná populace)}
Řetězce ztotožníme s~čísly $0..2^l-1$. Populaci popisuje vektor relativních četností $p(t)$
v~simplexu, vektor \emph{selekčních pravděpodobností} $s(t)$.
Hledáme $\mathcal{G}=\mathcal{M}\circ\mathcal{F}$ s~$p(t+1)=\mathcal{G}(p(t))$:
\begin{itemize}
\item $\mathcal{F}$ (selekce) je dána diagonální maticí fitness,
  $s=\mathbf{F}p/\sum_j \mathbf{F}_{jj}p_j$;
\item $\mathcal{M}$ (křížení $+$ mutace) je kvadratický operátor
  $\mathbb{E}[p'_k]=\sum_{i,j}s_is_j\,r(i,j,k)$; díky symetrii
  $r(i,j,k)=r(i\oplus k,j\oplus k,0)$ stačí \emph{jediná} míchací matice.
\end{itemize}
\tm{Pevné body}: $\mathcal{F}$ --- populace kopií nejlepšího jedince; $\mathcal{M}$ ---
rovnoměrné rozdělení. Jejich napětí vysvětluje \emph{přerušované rovnováhy}; pevné body celého
$\mathcal{G}$ známy nejsou.

\hh{Konečné populace: Markovův řetězec}
Stavy $=$ možné populace, jejich počet $N=\binom{n+2^l-1}{2^l-1}$; přechodová matice
$\mathbf{Q}$ je \emph{exaktní}, ale astronomicky velká.
Při $p_M>0$ je řetězec \tm{ireducibilní} (z~každé populace lze dosáhnout každé), takže optimum
je navštíveno s~pravděpodobností~1 --- ale bez elitismu se zase ztratí;
\tm{elitistický GA proto konverguje k~optimu s~pravděpodobností~1}.
Shrnutí: nekonečné populace $=$ deterministický model (exaktní, nepraktický), konečné $=$
stochastický (realistický, nezvládnutelný); limitní přechod je spojuje.

\hh{EDA --- praktický důsledek}
Algoritmy odhadu rozdělení berou pohled \uv{GA $=$ transformace rozdělení} doslova: místo
operátorů se \emph{učí model} $p_\theta$ a~z~něj vzorkují. Cyklus \emph{vzorkuj \ra{} vyber
lepší \ra{} přeuč model}.
\tm{UMDA} --- nezávislé marginály přepočítané z~vybrané podpopulace; \tm{PBIL} ---
inkrementálně $p_i\gets(1-\alpha)p_i+\alpha x_i^{\mathrm{best}}$; \tm{MIMIC/BMDA} --- párové
závislosti; \tm{BOA} --- bayesovská síť; \tm{CMA-ES} --- gaussovské rozdělení (otázka~3).

\end{multicols}
\newpage

%=====================================================================
\qpage{3}{Evoluční strategie, diferenciální evoluce, koevoluce, otevřená evoluce}
%=====================================================================
\begin{multicols}{2}
\footnotesize

\hh{Osnova odpovědi}
\ok{1} ES: charakteristické volby, notace, samoadaptace \ok{2} CMA-ES \ok{3} DE a~její
geometrická podstata \ok{4} koevoluce: typy, patologie, léky \ok{5} otevřená evoluce
a~novelty search.

\hh{Evoluční strategie}
Reálné vektory, \emph{mutace je hlavní operátor}, rodičovská selekce \emph{náhodná},
environmentální selekce \emph{deterministická}; jedinec $=[\vec x,\vec\sigma]$.
$(\mu,\lambda)$ vs.\ $(\mu+\lambda)$ --- umět rozdíly v~elitismu, konvergenci, schopnosti
opustit lokální optimum a~v~samoadaptaci (čárková varianta zapomíná, a~proto v~ní samoadaptace
funguje lépe: zastaralé $\sigma$ nepřežije).
\tm{$(1+1)$-ES a~pravidlo 1/5}: podíl úspěšných mutací $p>1/5$ \ra{} zvětši $\sigma$ (o~cca
10\,\%), $p<1/5$ \ra{} zmenši --- jde o~\emph{adaptivní}, nikoli samoadaptivní řízení.
\tm{Samoadaptace}: $\sigma'=\sigma\exp(\tau'N(0,1)+\tau N_i(0,1))$, teprve poté
$x'=x+\sigma'N(0,1)$ --- \emph{nejdřív $\sigma$, pak $x$}, aby se nový krok hned otestoval.
Počet strategických parametrů: 1 (izotropní), $n$ (nekorelované), $n(n+1)/2$ (korelované,
včetně úhlů natočení).
\tm{Rekombinace}: lokální ($\rho=2$) / globální ($\rho=\mu$), diskrétní (dominantní) /
intermediární (aritmetická).

\hh{CMA-ES}
Stav $\theta=(\vec m,\sigma,\mathbf C,\vec p_\sigma,\vec p_c)$, vzorkuje se
$\mathcal N(\vec m,\sigma^2\mathbf C)$.
Aktualizace: nový střed z~$\mu$ nejlepších vzorků; \tm{evoluční cesty} $\vec p_\sigma,\vec p_c$
akumulují posuny středu; \tm{rank-1} update podle $\vec p_c\vec p_c^{\top}$ a~\tm{rank-$\mu$}
update podle rozptylu vybraných jedinců; \emph{délka} $\vec p_\sigma$ řídí velikost kroku
$\sigma$.
Invariantní vůči monotónním transformacím fitness a~lineárním transformacím prostoru; de facto
standard pro spojitou black-box optimalizaci.

\hh{Diferenciální evoluce}
\tm{DE/rand/1/bin}: dárce $\vec v=\vec x_a+F(\vec x_b-\vec x_c)$; \emph{binomické} křížení
s~prahem $C$ a~pojistkou $j=I_{\mathrm{rand}}$ (alespoň jedna složka od dárce); souboj potomka
\emph{s~vlastním rodičem} --- nahrazení jen při zlepšení.
Parametry: $F\in[0,2]$ (typicky $0{,}5$--$0{,}9$), $C\in[0,1]$, $N\approx 10D$.
\tm{Geometrická podstata}: měřítko \emph{i} orientace kroků se přebírají z~aktuálního rozptylu
populace \ra{} adaptace bez strategických parametrů (široká populace = velké kroky,
zkonvergovaná = jemné doladění).
Varianty: rand/1, rand/2, best/1, best/2, current-to-best; \texttt{bin} vs.\ \texttt{exp}
křížení. \emph{Umět předvést krok s~čísly: mutace \ra{} křížení \ra{} selekce.}

\hh{Koevoluce}
$=$ kontextová fitness (hodnotí se \emph{proti ostatním}) $+$ dynamická krajina.
\tm{Kompetitivní}: predátor--kořist, host--parazit, Hillisovy řadicí sítě.
\tm{Kooperativní}: dekompozice na komponenty, CCGA, moduly $+$ blueprinty.
\tm{Patologie}: cyklení (intranzitivita), odpojení (jedna populace uteče), efekt
\emph{Červené královny} (všichni se zlepšují, ale \emph{relativní} fitness zůstává konstantní
--- z~naměřených hodnot nepoznáme, zda je pokrok skutečný), průměrné stabilní stavy,
přespecializace, zapomínání.
\tm{Léky}: archivy (\en{hall of fame}), sdílená fitness, měření proti \emph{pevné externí sadě}
protivníků.
\tm{Vězňovo dilema} jako modelová úloha: $T>R>P>S$ a~$2R>T+S$; $D$ dominantní, $DD$ je Nashovo
ekvilibrium a~jediné Pareto-neoptimální pole; iterovaná verze \ra{} TFT a~čtyři vlastnosti
úspěšné strategie (slušnost, dráždivost, odpouštění, nezávistivost) $+$ srozumitelnost jako
pátá. V~\emph{zašuměném} prostředí vítězí \tm{Pavlov} (win-stay, lose-shift: při výplatě $R$
nebo $P$ hraj $C$, při $T$ nebo $S$ hraj $D$).
Opakování Axelrodova turnaje po 20~letech vyhrál \emph{tým} ze Southamptonu: strategie se
prvními cca 10~tahy rozpoznaly a~pak všechny obětavě nahrávaly jedné z~nich.

\hh{Diverzita a~dynamické krajiny}
Diverzita: fitness sharing, crowding, speciace, ostrovní model, nenáhodné páření.
Dynamická fitness: diploidie, hypermutace, náhodná migrace, udržování diverzity,
$(\mu,\lambda)$-ES (zapomínání je zde výhodou).

\hh{Otevřená evoluce}
Evoluce \emph{bez pevného cíle}. Znaky: trvalá novost, neomezený růst složitosti, vznik nových
tříd entit, evoluce evolvability. Systémy Tierra, Avida, Polyworld, ECHO --- všechny se dříve
či později \emph{zaseknou}.
\tm{Generativní (nepřímé) reprezentace}: L-systémy, celulární automaty, buněčné kódování, CPPN.
\tm{Novelty search}: účelová funkce se \emph{vůbec nepoužívá}, měří se řídkost $\rho(x)$
v~prostoru \emph{chování} vůči populaci \emph{a} permanentnímu archivu --- u~deceptivních úloh
se vyplatí cíl zahodit. Navazuje \en{quality-diversity} a~\tm{MAP-Elites} (mřížka deskriptorů
chování, v~každé buňce nejlepší jedinec).
\tm{Interaktivní evoluce}: fitness $=$ člověk (Biomorphs, Picbreeder, EndlessForms); omezením
je malá populace a~únava uživatele.

\end{multicols}
\newpage

%=====================================================================
\qpage{4}{Rojové optimalizační algoritmy}
%=====================================================================
\begin{multicols}{2}
\footnotesize

\hh{Osnova odpovědi}
\ok{1} co dělá roj rojem \ok{2} PSO: rovnice, parametry, topologie \ok{3} čím se PSO liší od GA
\ok{4} ACO: rovnice a~konstrukce řešení \ok{5} varianty ACO \ok{6} kdy který algoritmus.

\hh{Východiska rojové inteligence}
Mnoho jednoduchých agentů bez centrálního řízení; z~lokálních interakcí a~úprav prostředí
vzniká globálně rozumné chování. Klíčová slova: \tm{stigmergie} (komunikace přes prostředí),
\tm{emergence}, \tm{decentralizace}, \tm{kladná zpětná vazba} (feromon) tlumená \tm{odpařováním}
jako zápornou zpětnou vazbou. Přímá vazba na reaktivní architektury agentů (subsumpce, Steelsův
průzkumník Marsu) --- rojové algoritmy jsou MAS řešící optimalizaci.

\hh{PSO --- rovnice}
$$\vec v \gets w\vec v + c_1r_1(\vec p-\vec x) + c_2r_2(\vec g-\vec x),\qquad
\vec x \gets \vec x+\vec v$$
$w$ $=$ \tm{setrvačnost} (klesá $0{,}9\to0{,}4$: nejdřív explorace, pak exploatace);
$c_1(\vec p-\vec x)$ $=$ \tm{kognitivní} složka (osobní nejlepší pozice);
$c_2(\vec g-\vec x)$ $=$ \tm{sociální} složka (nejlepší v~okolí); typicky $c_1=c_2=2$,
$r_1,r_2\sim U(0,1)$ po složkách.
Divergenci brání omezení $v_{\max}$ nebo \tm{constriction} $\chi\approx0{,}729$.
\tm{Topologie}: \en{gbest} (úplný graf --- rychlá konvergence, riziko uvíznutí) vs.\ \en{lbest}
(kruh, mřížka --- pomalejší šíření informace, lepší explorace).

\hh{PSO vs.\ GA}
Žádné křížení ani mutace; částice mají \emph{paměť} ($\vec p_i$, $\vec g$); informace teče
\emph{jednosměrně} od lepších k~horším; populace se nenahrazuje, jen se pohybuje.

\hh{ACO --- rovnice}
Úloha se převede na hledání nejkratší cesty v~ohodnoceném grafu; mravenci řešení
\emph{konstruují} po hranách, nikoli vylepšují.
Pravděpodobnost přechodu $p_{ij}\propto\tau_{ij}^{\alpha}\,\eta_{ij}^{\beta}$, kde $\tau$ $=$
feromon (naučená zkušenost) a~$\eta=1/d_{ij}$ $=$ heuristika (viditelnost); $\alpha,\beta$ váží
zkušenost proti heuristice.
Aktualizace $\tau\gets(1-\rho)\tau+\sum_k Q/L_k$ --- \emph{odpařování} $+$ přídavek úměrný
kvalitě řešení (kratší cesta \ra{} víc feromonu).

\hh{Varianty ACO}
\tm{AS} (původní, Dorigo 1991: feromon aktualizují \emph{všichni} mravenci) $\cdot$
\tm{elitní} a~\tm{rank-based} (aktualizuje jen globálně nejlepší, resp.\ $N$ nejlepších s~vahou
podle pořadí) $\cdot$
\tm{ACS} (inspirace Q-learningem: globální update jen nejlepším; navíc \emph{lokální} update
při každém průchodu hranou, který feromon \emph{snižuje}, aby ostatní zkoušeli jiné cesty;
pseudonáhodné pravidlo --- s~pravděpodobností $q_0$ rovnou nejlepší hrana, jinak vzorkování
jako v~AS) $\cdot$
\tm{MAX--MIN} (meze $\tau_{\min},\tau_{\max}$ proti předčasné konvergenci).
\tm{Akce démona} $=$ centralizovaný krok mimo mravence, typicky lokální prohledávání nad
zkonstruovanými řešeními.
Pro spojité úlohy se prostor diskretizuje na podintervaly a~feromon navádí do lepších z~nich.

\hh{Další rojové algoritmy (nadstavba nad slidy EVA)}
\tm{ABC} (\en{artificial bee colony}): \emph{dělnice} u~zdrojů, \emph{pozorovatelky} vybírají
zdroj ruletou podle kvality, \emph{průzkumnice} po vyčerpání limitu neúspěšných pokusů zdroj
opustí a~restartují jej náhodně.
\tm{Firefly}: přitažlivost klesá exponenciálně se vzdáleností, takže se roj samovolně dělí na
podskupiny okolo lokálních optim --- vhodné pro multimodální úlohy.

\hh{Kdy který}
Spojitá úloha bez struktury \ra{} PSO nebo CMA-ES; konstrukce cesty či permutace, kde existuje
heuristika \ra{} ACO; \emph{dynamické} prostředí (mění se hrany, vypadávají uzly) \ra{} ACO,
protože odpařování samo zapomíná neplatnou zkušenost; multimodální krajina \ra{} firefly,
lbest PSO.
Praxe ACO: TSP, směrování vozidel a~paketů, skládání proteinů, detekce hran v~obraze.

\hh{Meze rojových metod}
Mnoho parametrů bez teorie, jak je nastavit; konvergenční záruky jen pro speciální případy
(u~PSO za omezujících předpokladů na $w,c_1,c_2$); u~PSO hrozí \emph{stagnace} roje, když
všechny částice splynou s~$\vec g$ (léky: restart, náboj/odpuzování, lbest topologie);
ACO je drahé na vyhodnocení, protože každý mravenec staví celé řešení.
\tm{Diskrétní PSO}: rychlost se interpretuje jako pravděpodobnost překlopení bitu
(sigmoida), případně jako posloupnost prohození u~permutací.

\ph{Pasti}
PSO nemá selekci --- populace se nikdy nenahrazuje, jen se pohybuje.
ACO řešení \emph{staví} po hranách, nevylepšuje hotové.
Lokální update v~ACS feromon \emph{snižuje} --- je to nástroj explorace, ne exploatace.
Stigmergie není komunikace zprávou, ale \emph{úpravou prostředí}.

\end{multicols}
\newpage

%=====================================================================
\qpage{5}{Memetické algoritmy, hill climbing, simulované žíhání}
%=====================================================================
\begin{multicols}{2}
\footnotesize

\hh{Osnova odpovědi}
\ok{1} hill climbing a~jeho tři nepřátelé \ok{2} simulované žíhání \ok{3} tabu search
\ok{4} scatter search \ok{5} memetické algoritmy: Lamarck vs.\ Baldwin \ok{6} ladění vs.\ řízení
parametrů \ok{7} kam to celé patří.

\hh{Hill climbing}
Prohledávání okolí jednoho řešení: \tm{steepest ascent} (nejlepší soused) $\cdot$
\tm{first improvement} (první lepší --- levnější) $\cdot$ \tm{stochastic} (soused
s~pravděpodobností úměrnou zlepšení) $\cdot$ \tm{s~náhodnými restarty}.
Tři nepřátelé: \emph{lokální optima}, \emph{plošiny} (žádný gradient, jen náhodná chůze),
\emph{hřebeny} (zlepšení jen kombinací více kroků najednou).

\hh{Simulované žíhání}
\tm{Metropolisovo kritérium}: horší soused se přijme s~pravděpodobností
$P=\min\bigl(1,e^{-\Delta/T}\bigr)$ --- při vysoké teplotě téměř náhodná chůze, při nízké téměř
hill climbing.
\tm{Rozvrh chlazení}: geometrický $T\gets\alpha T$ (v~praxi), logaritmický $T_k\propto1/\log k$
(nutný pro důkaz konvergence, prakticky nepoužitelný).
Pro \emph{pevné} $T$ konverguje rozdělení stavů k~Boltzmannovu $\propto e^{-f/T}$.

\hh{Tabu search (Glover 1986)}
Vždy se přejde do \emph{nejlepšího nezakázaného} souseda, i~když je horší než aktuální řešení
--- odtud schopnost opustit lokální optimum bez náhody.
\tm{Tabu list} zakazuje nedávné \emph{tahy} (swap měst, překlopení bitu, přidání hrany), ne celá
řešení --- je to úspornější a~zobecňuje to lépe. \tm{Tenure} $=$ počet iterací, po které tah
zůstává zakázán (statické $T$, nebo náhodné z~rozsahu).
\tm{Aspirační kritérium} zákaz zruší, vede-li tah k~novému nejlepšímu řešení.
Tři paměti: \tm{krátkodobá} (tabu list) $\cdot$ \tm{střednědobá} (intenzifikace nad elitními
řešeními) $\cdot$ \tm{dlouhodobá frekvenční} (diverzifikace --- často navštěvované komponenty
se penalizují nebo se od nich restartuje).

\hh{Scatter search}
Referenční množina se vybírá podle \emph{kvality i~diverzity}, nová řešení vznikají
\emph{aritmetickými kombinacemi} podmnožin a~následným vylepšením lokálním prohledáváním;
deterministické tam, kde EA používá náhodu.

\hh{Memetické algoritmy (Moscato 1989)}
$=$ EA $+$ lokální prohledávání nad jedinci: populace hledá globálně, \emph{mem} (hill climbing,
SA, tabu, 2-opt, gradient/backpropagation) dolaďuje lokálně. Genetické operátory navíc
\emph{přenášejí informaci mezi běhy lokálního hledání}, což je celý zisk. Dvě varianty:
\begin{itemize}
\item \tm{lamarckismus} --- vylepšený jedinec \emph{přepíše genotyp}; rychlejší konvergence, ale
  rychlejší ztráta diverzity (a darwinisticky \uv{nekošer});
\item \tm{baldwinismus} --- vylepšení se projeví \emph{jen ve fitness}, genotyp zůstane;
  pomalejší, zachovává diverzitu a~vyhlazuje krajinu (fitness $=$ \emph{potenciál} jedince).
\end{itemize}
Rozhodnutí návrhu: na koho lokální hledání pustit (všichni / nejlepší / náhodně
s~pravděpodobností $p$), jak dlouho $t$ a~jak silně.
\tm{Memetic computing}: memy nejsou známy dopředu --- \emph{multi-meme MA} (mem je součást
genotypu a~dědí se) a~\emph{meta-Lamarckian MA} (memy mají vlastní populaci a~koevolvují
s~jedinci, odměnou memu je fitness jedince, který vylepšil).

\hh{Ladění vs.\ řízení parametrů}
\tm{Ladění} (\en{tuning}) --- hledá se \emph{před} během (grid search, iterované závody,
meta-EA); parametry ale nejsou nezávislé a~vyčerpávající hledání je nereálné.
\tm{Řízení} (\en{control}) --- mění se \emph{za} běhu:
\begin{itemize}
\item \tm{pevné} --- deterministicky podle času ($T$ v~SA, $w$ v~PSO, $\sigma(t)=1-0{,}9\,t/T$),
  případně stochasticky;
\item \tm{adaptivní} --- podle zpětné vazby z~běhu (pravidlo 1/5, Boltzmannův turnaj
  s~teplotou, kterou lze při uvíznutí zase zvýšit);
\item \tm{samoadaptivní} --- parametr je součástí genomu a~evolvuje ($\sigma$ v~ES).
\end{itemize}
Řídit lze i~\emph{reprezentaci} (přidávat bity, klesne-li hammingovská diverzita) a~\emph{velikost
populace} (jedinec dostane délku života úměrnou fitness a~pak zemře).

\hh{Souvislosti}
SA, tabu a~HC jsou \emph{jednobodové} metody, které EA doplňují jako lokální krok (memetické
algoritmy, akce démona v~ACS, 2-opt v~TSP, WalkSAT v~SAT). U~kombinatorických úloh je
hybridizace s~lokální heuristikou téměř vždy výhodná (otázka~6).

\end{multicols}
\newpage

%=====================================================================
\qpage{6}{Aplikace evolučních algoritmů (evoluce expertních systémů, neuroevoluce,
řešení kombinatorických úloh, vícekriteriální optimalizace)}
%=====================================================================
\begin{multicols}{2}
\footnotesize

\hh{Osnova odpovědi}
\ok{1} LCS: Michigan vs.\ Pittsburgh, ZCS, XCS \ok{2} neuroevoluce: co se evolvuje, konkurující
konvence, NEAT a~HyperNEAT \ok{3} kombinatorické úlohy: rozhoduje kódování $+$ práce s~omezeními
\ok{4} vícekriteriální optimalizace: dominance, skalarizace, NSGA-II.

\hh{LCS: dva přístupy}
\tm{Michigan}: jedinec $=$ \emph{jedno pravidlo}, řešením je celá populace \ra{} nutné rozdělení
odměny (\en{credit assignment}).
\tm{Pittsburgh}: jedinec $=$ \emph{celá sada pravidel}, fitness přímočará, ale operátory složité
(GIL).
\tm{Hollandův LCS}: podmínky nad $\{0,1,*\}$, vnitřní zprávy a~receptory, \tm{bucket brigade}
(pravidlo platí předchůdci za právo jednat).
\tm{ZCS}: bez zpráv; \en{match set} $M$ \ra{} ruletový výběr akce \ra{} \en{action set} $A$;
posílení $A$ a~$A_{-1}$, daň pro $M\setminus A$; operátor \en{cover} při prázdném $M$.
Slabiny: neevolvuje \emph{úplnou} sadu pravidel, pravidla na začátku řetězců se odměňují zřídka
a~vymírají, stejně tak pravidla vedoucí k~malým, ale nutným odměnám.
\tm{XCS}: klasifikátor $(c,a,p,\epsilon,F)$, \emph{fitness podle přesnosti predikce}
$\kappa=\alpha(\epsilon/\epsilon_0)^{-\nu}$ (pro $\epsilon\le\epsilon_0$ je $\kappa=1$; fitness
je \emph{relativní} přesnost v~action setu), predikce se aktualizuje Q-learningem, GA běží
\emph{nikově} uvnitř action setu.
Výsledkem je \emph{úplná a~minimální} mapa vstup--výstup --- to je hlavní pointa oproti
\uv{strength-based} systémům, kde silné, ale nepřesné pravidlo vyhrávalo.
Varianty UCS (s~učitelem), XCSF (aproximace funkcí).

\hh{Neuroevoluce}
Co se evolvuje: váhy / topologie / obojí / hyperparametry učení. Přednosti EA: nepotřebuje
gradient, zvládá rekurentní sítě a~RL s~řídkou odměnou, snadno se paralelizuje.
\tm{Problém konkurujících konvencí} (permutační problém): dvě funkčně stejné sítě mají různě
uspořádané genomy \ra{} křížení je bez opatření destruktivní.
\tm{Kódování architektury}: přímé (matice vah), gramatické (Kitano), buněčné (Gruau --- GP
program řídí růst sítě); \tm{GNARL} $=$ EP nad grafy, strukturální i~parametrické mutace řízené
\uv{teplotou}.
\tm{NEAT}: seznam hran s~\emph{inovačními čísly} (křížení jen shodných genů --- tím řeší
konkurující konvence), \emph{speciace} podle podobnosti genomů se sdílenou fitness (chrání nové
topologie, než se stihnou vyladit), \emph{complexification} od minimální sítě.
\tm{HyperNEAT}: neurony leží v~geometrickém \emph{substrátu} a~váhu spoje generuje z~jejich
souřadnic funkce $\mathbb{R}^4\to\mathbb{R}$ reprezentovaná \tm{CPPN} (DAG s~různými aktivačními
funkcemi) evolvovanou pomocí NEAT --- zachytí symetrie a~regularity, oddělí velikost genomu od
velikosti sítě.
\tm{NAS}: prostor (lineární / DAG buněk / super síť), strategie (EA, RL, bayesovská, gradientní),
odhad výkonu (sdílení vah --- ENAS); DeepNEAT a~CoDeepNEAT (moduly $+$ blueprinty); ES jako
škálovatelná alternativa gradientu v~deep RL.

\hh{Kombinatorické úlohy}
\tm{Batoh}: bitová mapa, triviální operátory, problém je \emph{kapacitní omezení} --- čtyři
strategie: \tm{penalizace} (volba váhy $\alpha$), \tm{oprava} (lamarckovsky nebo baldwinovsky),
\tm{dekodér} (\uv{1 znamená: přidej, pokud se vejde} --- \emph{každý} genotyp je pak přípustný),
\tm{vícekriteriální} formulace (maximalizuj cenu, minimalizuj překročení).
\tm{TSP}: fitness triviální, rozhoduje kódování. Sousednostní a~ordinální reprezentaci
nepoužívat; permutační $+$ PMX/OX/CX/ER, případně maticová. Inicializace nejbližším sousedem či
vkládáním hran; mutace inverze (2-opt).
\tm{SAT}: fitness $=$ počet splněných klauzulí (počet splněných \emph{proměnných} nestačí ---
plošiny), dále vážení klauzulí a~\emph{zjemňující funkce} $f+\alpha r(x)$, která rozliší jedince
se stejnou fitness (váhy $v_j$ sledují tendenci proměnné k~0/1). Reprezentace bitová, reálná,
klauzulová, cestová; WalkSAT jako lokální heuristika. U~reálné reprezentace
($x\mapsto(1-y)^2$, $\neg x\mapsto(1+y)^2$, minimalizace) pozor: hodnoty jsou \emph{penalty},
takže disjunkce (klauzule) je \tm{součin} a~konjunkce (formule) \tm{součet}.
\tm{Rozvrhování}: přímé maticové kódování; \emph{tvrdá} omezení řešit v~operátorech,
\emph{měkká} ve fitness; job shop --- permutace s~dekodérem vs.\ přímé kódování celého plánu.
Obecné pravidlo: rozhoduje kódování a~operátory respektující strukturu úlohy $+$ hybridizace
s~lokální heuristikou.

\hh{Vícekriteriální optimalizace}
Místo jedné fitness vektor $f_1,\dots,f_n$ (uvažujme minimalizaci).
\tm{Slabá dominance} $f_i(x)\le f_i(y)$ pro všechna $i$; \tm{dominance} navíc
$\exists j: f_j(x)<f_j(y)$; jinak jsou řešení \tm{neporovnatelná}.
\tm{Pareto množina} (v~prostoru řešení) a~\tm{Pareto fronta} (v~prostoru kritérií, množina
nedominovaných). Dvojí cíl: \emph{konvergence} k~frontě $+$ \emph{rovnoměrné pokrytí} fronty.
\tm{Skalarizace}: lineární vážený součet (jednoduchá, ale nenajde nekonvexní části fronty
a~váhy neumíme zvolit), \emph{$\varepsilon$-constraint} (minimalizuj $f_1$, ostatní jako
omezení), Čebyševova.
\tm{VEGA} (1985, první): populace se rozdělí na $n$ dílů po $N/n$ nejlepších podle každého
$f_i$ --- konverguje k~jednotlivým optimům a~ztrácí diverzitu.
\tm{NSGA} (1994): nedominované třídění do front $F_1,F_2,\dots$ ($F_1$ nedominovaní z~$P$,
$F_2$ z~$P\setminus F_1$, \dots) $+$ nikový faktor $\mathrm{sh}(i,j)=1-(d_{ij}/d_{\text{share}})^2$
pro $d_{ij}<d_{\text{share}}$; slabiny: volba $d_{\text{share}}$ a~chybějící elitismus.
\tm{NSGA-II} (2000) obojí opravuje: místo $d_{\text{share}}$ \emph{crowding distance}
$d_i=\sum_m\bigl(f_m(i{+}1)-f_m(i{-}1)\bigr)/(f_m^{\max}-f_m^{\min})$, krajní body $\infty$;
\emph{crowded comparison} (nižší fronta vyhrává, při shodě větší crowding) v~turnaji --- žádná
fitness se nepočítá, porovnávají se jen ta dvě čísla; elitismus spojením $P_t\cup Q_t$.
\emph{Umět spočítat crowding distance na příkladu.}
\tm{NSGA-III}: pro \en{many-objective} nahrazuje crowding \emph{referenčními body}
(rovnoměrně rozmístěné směry, počet $\approx$ velikost populace, $N=\binom{p+M+1}{p}$);
při plnění se přednostně doplňují nedostatečně zastoupené směry.
\tm{MOEA/D}: dekompozice na podpopulace s~různými váhovými vektory $\lambda^{(i)}$, Čebyševova
skalarizace $\min\max_i\lambda_i|f_i(x)-z_i|$ vůči referenčnímu bodu $z$, křížení jen
v~\emph{sousedství} $B(i)$ nejbližších váhových vektorů.
\tm{SMS-EMOA}: steady-state, z~$P\cup\{q\}$ se vybere podmnožina velikosti $\mu$ maximalizující
\tm{hypervolume} (S-metriku). \tm{SPEA2}: archiv, síla dominujících, hustota přes $k$-tého
souseda.
\tm{Hypervolume} $=$ objem dominovaný frontou vůči referenčnímu bodu --- jediná běžná metrika
monotónní vůči dominanci a~nepotřebující znát skutečnou frontu; dále GD, IGD, spread,
$\varepsilon$-indikátor.

\end{multicols}

\end{document}
